\documentclass[11pt]{article}

\usepackage[margin=1in]{geometry}
\usepackage{amsmath,amssymb}
\usepackage{booktabs}
\usepackage[T1]{fontenc}
\usepackage{lmodern}
\usepackage{microtype}
\usepackage{hyperref}
\usepackage{enumitem}
\usepackage{listings}
\usepackage{xcolor}

\hypersetup{%
  colorlinks=true,
  linkcolor=blue,
  urlcolor=blue,
  citecolor=blue
}

\lstdefinestyle{py}{%
  language=Python,
  basicstyle=\ttfamily\small,
  columns=fullflexible,
  breaklines=true,
  frame=single,
  rulecolor=\color{black!20},
  keywordstyle=\color{blue!60!black},
  commentstyle=\color{black!60},
  stringstyle=\color{green!40!black}
}

\title{FieldSemantics From First Principles}
\author{\emph{(Seth Price)}}
\date{\today}

\setlength{\parindent}{0pt}
\setlength{\parskip}{0.6\baselineskip}

\begin{document}
\maketitle

\begin{abstract}
FieldSemantics is a local, operational account of how ``meaning'' is produced under constraints: in frames, with budgets, through transforms. Each construct below comes with an explicit handle: you can run the operators, audit the costs, and predict failure modes.
\end{abstract}

\section{Running example and complexity ladder}

We use one artifact and progressively increase layer depth.

\textbf{Artifact $Z$, Layer 0:} a single requirement line in a product document:

\begin{quote}
\textbf{$Z_0$:} ``The app must work offline.''
\end{quote}

\textbf{Why this example?}
\begin{itemize}[leftmargin=2em]
  \item It is short (high compression already).
  \item It crosses stakeholders (multiple apertures).
  \item It forces integrability questions (can all frames stitch to one implementation?).
  \item It reliably produces ``demons'' (systems that keep ``alignment'' stable by exporting cost).
\end{itemize}

We add layers as we go:
\begin{itemize}[leftmargin=2em]
  \item Layer 0: the sentence alone.
  \item Layer 1: frame indexing (who reads it, when).
  \item Layer 2: meaning as a vector (not a scalar).
  \item Layer 3: aperture and orientation (what is admitted, what gradients are tracked).
  \item Layer 4: map/reduce (branching and stabilization).
  \item Layer 5: calculus (integrability, curvature, path dependence).
  \item Layer 6: orientability and loop-transport.
  \item Layer 7: demons (invariant maintenance by cost export).
  \item Layer 8: trace vs.~image (why artifacts lie even when everyone is sincere).
\end{itemize}

\section{Conventions: five nouns you must keep distinct}

Before doing anything ``semantic,'' we fix types. If you blur these, you will hallucinate global meaning.

\begin{description}[leftmargin=2.2em,labelsep=0.6em]
  \item[Entity:] a stable referent used for operations (not essence). \\
  Example: ``offline mode'' as a handle, not a metaphysical object.

  \item[Operator:] a transform between states/structures. \\
  Example: ``turn this requirement into a ticket.''

  \item[State:] an instantaneous configuration at a frame. \\
  Example: an engineer's current interpretation at 10:12 AM.

  \item[Trace:] a sequence of states/operators over frames. \\
  Example: ``PM wrote $\rightarrow$ ENG interpreted $\rightarrow$ QA tested $\rightarrow$ Support triaged.''

  \item[Cost:] distortion, hidden work, or denied budget. \\
  Example: ``QA discovers we did not define `offline'---now QA pays the definition cost under deadline.''
\end{description}

\textbf{If you learn nothing else:} meaning lives in traces, not in single statements.

\section{Notation: minimal alphabet, maximum leverage}

We use a compact set of symbols; every one has an operational role.

\begin{center}
\begin{tabular}{@{}ll@{}}
\toprule
Symbol & Definition \\
\midrule
$\sigma$ & Live in-frame semantic structure (the field configuration) \\
$\sigma_e$ & Externalized / frozen structure (exported artifact) \\
$F_t$ & A frame (a discrete evaluation slice) \\
$\Delta$ & Local differential move (map-step; branching or micro-update) \\
$\int$ & Stabilization across frames (reduce-step; integration into canonical form) \\
$A$ & Aperture (what evidence/axes are admitted or excluded) \\
$O$ & Orientation (tangent alignment; which gradients are tracked) \\
$C$ & Compression (lossy reduction) \\
$D$ & Distortion (mismatch induced by compression/transport/misorientation) \\
$K$ & Curvature (nonlinearity, thresholds, path dependence) \\
$S$ & Subject-position (what counts as ``given'' and ``actionable'') \\
$I$ & Image (the exported surface artifact people point at) \\
\bottomrule
\end{tabular}
\end{center}

You can read the whole project as: track $\sigma$, audit $A/O/C/D$, detect $K$, and refuse to confuse $I$ for the trace $T$ that produced it.

\section{Primitive 1: Locality}

\textbf{Claim:} meaning-operations are defined locally (frame-bounded) and must be measurable by observables inside that frame. A frame $F_t$ is a discrete slice: a moment plus the information available in that moment.

\subsection{Example upgrade: Layer 1}

Take $Z_0$: ``The app must work offline.'' Now index it by frames:
\begin{itemize}[leftmargin=2em]
  \item $F_0$ (PM writing): the product manager writes $Z_0$ at planning time.
  \item $F_1$ (ENG reading): engineering reads $Z_0$ in the sprint backlog.
  \item $F_2$ (QA reading): QA derives test cases from $Z_0$.
  \item $F_3$ (Marketing reading): marketing turns $Z_0$ into copy.
  \item $F_4$ (Support reading): support receives tickets: ``offline broken.''
\end{itemize}

Same string, different frames. FieldSemantics starts here: no statement has ``the meaning'' without a frame index.

\textbf{Firmware load:} always ask: ``What frame are we in, and what inputs are locally available?'' If someone answers without specifying a frame, they are doing global closure by default.

\section{Primitive 2: Vectorial semantics}

Most discourse tries to scalarize meaning: ``What does this mean?'' $\rightarrow$ one answer. FieldSemantics refuses. Meaning has components.

\subsection{Operational definition}

Let meaning in a frame be a vector $m \in \mathbb{R}^n$. Each component is a trackable axis (intent-channel, constraint-channel, action-channel, etc.). You choose axes based on the task, but you must name them.

\subsection{Example upgrade: Layer 2}

For $Z_0$ (``work offline''), define a minimal meaning vector:
\[
  m =
  \begin{bmatrix}
    u\\ v\\ w\\ p\\ r
  \end{bmatrix}
  \quad\equiv\quad
  \begin{bmatrix}
    \text{Network independence}\\
    \text{Freshness / staleness tolerance}\\
    \text{Auth + entitlement behavior}\\
    \text{Product coverage (which flows)}\\
    \text{Risk constraints (security/compliance)}
  \end{bmatrix}.
\]

Now you can express differences between frames as different vectors, not arguments about ``what it means.''

\begin{itemize}[leftmargin=2em]
  \item In $F_0$ (PM), $m$ might weight $u$ (``usable on planes'') and $p$ (``core flows work''), while underweighting $w,r$.
  \item In $F_1$ (ENG), $m$ might heavily weight $w$ (token caching, expiry edge cases) and $r$ (encryption, threat model).
  \item In $F_3$ (Marketing), ``offline'' may collapse into $u$ only (``no internet needed'') while silently deleting $v,w,r$.
\end{itemize}

\textbf{Firmware load:} ask: ``What are the axes of meaning here?'' If the axes are not named, you will fight over a scalar that never existed.

\section{Primitive 3: Aperture and Orientation}

\textbf{$A$ chooses what exists; $O$ chooses what matters.}

\subsection{Aperture $A$}

$A$ is a selection function: what evidence/axes are admitted as ``in play.'' Operationally, you can represent aperture as a set of channels, or as a projection operator $P_A$ that zeroes excluded dimensions. If $m\in\mathbb{R}^n$, an aperture can act like
\[
  m_A = P_A m.
\]

\subsection{Orientation $O$}

$O$ is which directions of change are treated as salient---what gradients the system follows. Operationally, orientation can be a weighting metric $G_O$ on meaning space, or an update rule that privileges some derivatives. For small changes,
\[
  \Delta m \propto -G_O\,\nabla \Phi(m),
\]
where $\Phi$ is a local objective or ``energy'' function (shipping time, risk exposure, user friction, cost, etc.).

\subsection{Example upgrade: Layer 3}

\textbf{$F_0$ (PM) aperture might be:}
\begin{itemize}[leftmargin=2em]
  \item Admits: user story, competitive checkbox, ``works on plane'' anecdote.
  \item Excludes: auth caching details, eventual consistency, entitlement rules.
\end{itemize}
So PM's $P_A$ effectively drops $w,r$ from the working meaning state.

\textbf{$F_1$ (ENG) aperture might be:}
\begin{itemize}[leftmargin=2em]
  \item Admits: token expiry, caching layers, storage limits, encryption requirements.
  \item Excludes: marketing promise tone, strategic positioning.
\end{itemize}
So ENG's $P_A$ drops rhetorical story dimensions and retains $w,r$.

\textbf{Orientation mismatch} happens when both admit overlapping axes but track different gradients. Example: both track ``core flows'' ($p$), but PM's gradient is ``minimize user friction,'' while ENG's gradient is ``minimize attack surface.'' Same axis, opposite update direction.

\textbf{Firmware load:}
\begin{enumerate}[leftmargin=2em]
  \item Ask: ``What is being excluded?'' ($A$)
  \item Ask: ``What changes trigger updates?'' ($O$)
\end{enumerate}
When two people ``disagree,'' it is often $A/O$ mismatch, not ignorance.

\section{Primitive 4: Compression and Distortion}

\textbf{Export requires compression; compression induces distortion.}

\subsection{Compression $C$}

$C$ is a lossy mapping from rich structure to reduced representation:
\[
  C:\sigma \rightarrow \rho,
\]
where $\rho$ is a reduced form (ticket, summary, label, OKR bullet).

\subsection{Distortion $D$}

$D$ is the mismatch you get when you use $\rho$ as if it were $\sigma$:
\[
  D = \operatorname{error}(\sigma,\ \operatorname{lift}(\rho)),
\]
where ``lift'' means reconstructing a usable structure from the reduced form.

\subsection{Example upgrade: Layer 4}

$Z_0$ becomes a backlog item:
\begin{quote}
  \textbf{$\sigma_e$:} ``Implement offline mode.''
\end{quote}

This is Freeze/Export (formalized next). Notice what happened:
\begin{enumerate}[leftmargin=2em]
  \item $C$ collapsed a multi-axis vector $m$ into a phrase.
  \item The omitted axes reappear later as bugs, scope fights, or customer confusion.
\end{enumerate}

Distortion here is operational, not philosophical:
\begin{itemize}[leftmargin=2em]
  \item QA writes tests assuming ``offline = all flows work.''
  \item ENG implements ``offline = cached read-only.''
  \item Marketing implies ``offline = fully functional.''
  \item Users hit the mismatch surface.
\end{itemize}

\textbf{Firmware load:} any time you see a short label, assume a compression operator ran. Then ask: ``What got dropped---and who will pay to reconstruct it?''

\section{Primitive 5: Map/Reduce}

\textbf{$\Delta$ branches locally; $\int$ stabilizes across frames.}

\subsection{Map step $\Delta$}

\[
  \operatorname{Map}:\sigma \rightarrow \{\sigma_i\}.
\]

A healthy Map step produces multiple candidate interpretations/implementations.
\begin{itemize}[leftmargin=2em]
  \item Failure mode: mode collapse (prematurely choosing one interpretation).
\end{itemize}

\subsection{Reduce step $\int$}

\[
  \operatorname{Reduce}:\{\sigma_i\} \rightarrow \sigma^* \quad \text{or} \quad \sigma_e.
\]

Reduce chooses a canonical form: a decision, a spec, a release note, a ``meaning.''
\begin{itemize}[leftmargin=2em]
  \item Failure mode: forced coherence (smoothing over contradictions without accounting for remainder).
\end{itemize}

\subsection{Example upgrade: Layer 5}

From $Z_0$, Map generates plausible implementations:
\begin{itemize}[leftmargin=2em]
  \item $\sigma_1$: offline = cached read-only content (no mutations).
  \item $\sigma_2$: offline = full CRUD with sync conflict resolution.
  \item $\sigma_3$: offline = ``no login required'' (marketing interpretation).
  \item $\sigma_4$: offline = ``works when network is flaky'' (degraded mode).
\end{itemize}

Then Reduce chooses one, e.g.~$\sigma_1$, and exports:
\begin{quote}
  ``Offline mode (read-only) shipping in v2.3''
\end{quote}

If Reduce does not encode the choice \emph{and the rejected branches}, recipients will reconstruct missing context incorrectly.

\textbf{Firmware load:} force a Map: ``List at least 3 $\sigma_i$.'' Then audit Reduce: ``Which $\sigma_i$ did we collapse to, and what remainder did we deny?''

\section{Freeze/Export as an operator}

\textbf{$\sigma \rightarrow \sigma_e$ is a real transform with predictable failure modes.}

We formalize what ``writing it down'' is doing:
\[
  \operatorname{Freeze}:\sigma \rightarrow \sigma_e.
\]

\begin{itemize}[leftmargin=2em]
  \item Local behavior ($\Delta$): choose a representational basis; apply compression $C$.
  \item Integral effect ($\int$): stabilize across frames (document persists).
  \item Failure modes: reification (treat $\sigma_e$ as $\sigma$), denial of distortion $D$.
\end{itemize}

\subsection{Example upgrade: Layer 6}

If the PRD chooses the basis ``feature checklist,'' offline becomes a checkbox. If it chooses ``user journey narrative,'' offline becomes story coverage. If it chooses ``threat model,'' offline becomes a security constraint problem.

Same ``offline,'' different basis $\rightarrow$ different compression $\rightarrow$ different distortion.

\textbf{Firmware load:} ask: ``What representational basis did we freeze into?'' If you cannot name the basis, you cannot predict the distortion surface.

\section{Calculus kit}

The goal is not to cosplay math. The goal is to make path dependence and stitch failures legible.

\subsection{Local differential}

A derivative is ``how meaning changes under a small change in inputs.'' Let $x$ be a context vector (constraints, examples, metadata). Meaning state $m(x)$ changes with $x$. The Jacobian is
\[
  J(x)=\frac{\partial m}{\partial x}.
\]
Operational translation: what features move which meaning axes?

\subsection{Gradient}

If you define a local objective (or energy) $\Phi(m)$, the gradient is $\nabla\Phi(m)$. Operational translation: which direction the system tries to move meaning to reduce local tension/cost.

\subsection{Integral}

An integral is accumulation: summing small updates across time/frames.
\[
  m_{t+1}=m_t + \Delta m_t \quad \Rightarrow \quad m_T = m_0 + \sum_{t=0}^{T-1}\Delta m_t.
\]
Operational translation: your ``final meaning'' is accumulated micro-updates under constraints.

\subsection{Integrability}

A field of updates is \emph{integrable} if local updates can be accumulated into a global potential in a path-independent way. In calculus terms: a vector field $F$ is integrable (conservative) on a region if there exists $\Phi$ such that $F=\nabla\Phi$, and therefore line integrals are path independent:
\[
  \int_{\gamma_1} F\cdot d\ell = \int_{\gamma_2} F\cdot d\ell \quad \text{(same endpoints)}.
\]

\textbf{Operational test for integrability:}
\begin{enumerate}[leftmargin=2em]
  \item Take the same evidence set.
  \item Incorporate it in different orders / negotiation paths.
  \item If the stabilized meaning differs materially $\Rightarrow$ \emph{non-integrable region}.
\end{enumerate}

\subsection{Example: path dependence in ``offline''}

Suppose two constraints are introduced:
\begin{itemize}[leftmargin=2em]
  \item $c_1$: ``Works on planes'' (pushes toward local caching).
  \item $c_2$: ``Always up to date'' (pushes toward live sync guarantees).
\end{itemize}

\textbf{Path A:} incorporate $c_1$ first. Engineering proposes cached read-only ($\sigma_1$). Then $c_2$ arrives and gets reinterpreted as ``up to date when online,'' not ``always.''

\textbf{Path B:} incorporate $c_2$ first. Engineering proposes always-online assumptions. Then $c_1$ arrives and ``offline'' gets reinterpreted as ``graceful error messaging,'' not true offline.

Same constraints, different incorporation path $\Rightarrow$ different $\sigma^*$ $\Rightarrow$ non-integrable.

\textbf{Firmware load:} integrability is not a vibe. It is a path-independence test.

\section{Curvature $K$}

\textbf{Where linear reasoning fails: thresholds and regime switches.}

Curvature here means: local geometry is nonlinear; small changes can cause disproportionate shifts; history matters.

\textbf{Operational proxies:}
\begin{itemize}[leftmargin=2em]
  \item sensitivity spikes,
  \item thresholds (``one more stakeholder and the meaning flips''),
  \item hysteresis (you cannot get back by reversing a step).
\end{itemize}

\subsection{Example upgrade: Layer 7}

At first, everyone treats offline as ``nice-to-have.'' Then legal adds: ``Offline data must be encrypted at rest.'' Suddenly the feasible set changes; implementation cost jumps; timeline shifts; ``offline'' becomes ``read-only'' to ship.

That is curvature: a local constraint changes the geometry of the solution space. If you are doing linear reasoning (``offline is just a feature''), you will be surprised by the jump. If you are tracking curvature, you predict it.

\textbf{Firmware load:} ask: ``Where are the thresholds? What constraint flips the regime?''

\section{Orientability}

\subsection{Orientability (definition)}

A space is \emph{orientable} if you can choose a consistent orientation (a consistent ``handedness'') across the whole region. Non-orientable spaces have the property that if you travel around a loop, your orientation flips.

\textbf{Operational orientability test:}
\begin{enumerate}[leftmargin=2em]
  \item Pick a local ``positive direction'' for an axis (what counts as ``more of $u$'').
  \item Transport it through frames/teams and back.
  \item If you return with the axis inverted, the semantic manifold is non-orientable under those mappings.
\end{enumerate}

\subsection{Semantic orientability}

In FieldSemantics, orientation $O$ was ``which gradients are tracked.'' Orientability asks a deeper question: \emph{can different frames share a consistent sign convention and update direction for the same axis?}

A useful way to formalize transport is to define mappings between bases $T_{i\rightarrow j}$, which map an axis convention from frame/team $i$ to $j$. A loop transport is
\[
  T_{\text{loop}} = T_{\text{PM}\rightarrow \text{Mkt}}\, T_{\text{Mkt}\rightarrow \text{Sup}}\, T_{\text{Sup}\rightarrow \text{PM}}.
\]

Orientable means $T_{\text{loop}}$ preserves axis sign conventions (approximately identity on relevant axes). Non-orientable means some axis comes back flipped.

\subsection{Example upgrade: Layer 8}

\textbf{Team A uses ``offline'' to mean:}
\begin{itemize}[leftmargin=2em]
  \item $(+u)$ = more network independence (good),
  \item shipping goal: increase $u$.
\end{itemize}

\textbf{Team B uses ``offline'' to mean:}
\begin{itemize}[leftmargin=2em]
  \item $(+w)$ = fewer login requirements / fewer entitlements checks (good),
  \item shipping goal: increase $w$.
\end{itemize}

These are not merely different emphases; they can invert each other in practice: increasing $u$ may require more auth caching complexity (which increases security surface), while increasing $w$ may require less enforcement (which increases risk).

Now run the loop PM $\rightarrow$ Marketing $\rightarrow$ Support $\rightarrow$ PM. By the time ``offline'' returns to PM via customer complaints, ``offline'' now means ``cannot login,'' not ``no network.'' The axis you thought was positive got swapped or inverted.

That is a semantic non-orientability event.

\textbf{Firmware load:} if your system is non-orientable, you cannot get global agreement by ``defining terms once.'' You need explicit transport rules (mappings) between frames.

\section{Demons}

\textbf{A demon is an operator that keeps invariants stable by exporting cost.}

A demon is not a moral category. It is a control primitive: a hidden operator that preserves a broken invariant by exporting cost elsewhere (or denying it).

\textbf{Operational signatures:}
\begin{itemize}[leftmargin=2em]
  \item an invariant claim persists despite contradictory evidence,
  \item displaced labor appears (usually downstream),
  \item opacity is part of the mechanism.
\end{itemize}

\subsection{Example upgrade: Layer 9}

\textbf{Invariant the org wants to preserve:} ``Requirements are clear.'' ``Teams are aligned.'' ``Scope is controlled.''

\textbf{Local reality:} ``offline'' means different things across frames. Reduce chose $\sigma_1$ but did not export the decision boundaries.

\textbf{The demon-maintain operator appears as:}
\begin{itemize}[leftmargin=2em]
  \item templated docs that look complete,
  \item meetings that end with ``we're aligned,''
  \item tickets that compress disagreements into one line.
\end{itemize}

\textbf{Cost export:}
\begin{itemize}[leftmargin=2em]
  \item QA pays with undefined edge cases.
  \item Support pays with angry users.
  \item Engineers pay with last-minute scope triage.
  \item PM pays with timeline slips---but only when the demon cannot hold anymore.
\end{itemize}

\textbf{Firmware load:} when you see a stable invariant in the presence of local contradiction, assume a demon and ask: ``Where is the cost being routed?''

\section{Subject-position $S$ and Image $I$}

$S$ (subject-position) is the stance that decides what is ``given'' and ``actionable.'' It sets aperture and orientation implicitly. $I$ (image) is the exported surface artifact: the PRD line, the ticket title, the release note.

\subsection{Example upgrade: Layer 10}

\textbf{Image $I$:} ``Offline mode shipped.''

\textbf{Trace $T$ (what actually happened):}
\begin{itemize}[leftmargin=2em]
  \item $F_0$: PM wrote $Z_0$ with axes $u,p$ salient.
  \item $F_1$: ENG mapped to $\sigma_1$ (cached read-only) due to security constraints.
  \item $F_2$: Reduce chose $\sigma_1$ but did not export disallowed $\sigma_2$ (full sync) as a rejected branch.
  \item $F_3$: Marketing copy implied $\sigma_2$.
  \item $F_4$: Support tickets forced a second Reduce under crisis, collapsing to ``offline = view-only'' publicly.
\end{itemize}

Two teams fight over ``what offline means'' because they only see $I$. FieldSemantics insists you carry $T$.

\textbf{Firmware load:} if you cannot reconstruct $T$, you will hallucinate causality from $I$.

\section{A reader-runnable procedure}

\textbf{Input:} an artifact $Z$ (sentence, label, policy line, meme, metric).

\begin{enumerate}[leftmargin=2em]
  \item Frame it: enumerate $F_t$ (who reads it, when, under what constraints?).
  \item Name meaning space: define axes ($m\in\mathbb{R}^n$).
  \item Declare aperture $A$: what channels/axes are admitted? excluded?
  \item Declare orientation $O$: what changes trigger updates? what gradients matter?
  \item Map ($\Delta$): generate $\{\sigma_i\}$: at least 3 plausible live structures.
  \item Reduce ($\int$): select $\sigma^*$. Explicitly record rejected branches and boundary conditions.
  \item Audit compression $C$: what got dropped in export to $\sigma_e$?
  \item Measure distortion $D$: where will mismatch surface? who will pay?
  \item Test integrability: reorder evidence / change negotiation path. If outcome differs $\Rightarrow$ non-integrable region.
  \item Test orientability: transport sign conventions around stakeholder loops. If axes invert $\Rightarrow$ non-orientable mapping; require explicit transport rules.
  \item Scan for demons: identify invariants held stable; locate displaced cost and opacity.
  \item Preserve trace: publish $T$ alongside $I$ when stakes are real.
\end{enumerate}

\section{Minimal working pseudocode}

\begin{lstlisting}[style=py]
def fieldsemantics_bootstrap(artifact_Z):
    # 1) Frames
    frames = enumerate_frames(artifact_Z)  # F_t

    # 2) Meaning space (vectorial)
    axes = choose_axes(artifact_Z)         # m in R^n

    # 3) Aperture & orientation per frame
    models = []
    for F in frames:
        A = set_aperture(F, axes)
        O = set_orientation(F, A)
        σ = infer_live_structure(F, artifact_Z, A, O)
        models.append((F, A, O, σ))

    # 4) Map: branch generation
    branches = map_branches(models)        # {σ_i}

    # 5) Reduce: stabilization + explicit remainder
    σ_star, remainder = reduce_with_remainder(branches)

    # 6) Export: σ -> σ_e
    σ_e = freeze_export(σ_star)

    # 7) Costs: compression, distortion, curvature, demons
    C = estimate_compression(σ_star, σ_e)
    D = estimate_distortion(σ_star, σ_e)
    K = detect_curvature(branches, remainder)
    demons = detect_demons(frames, σ_e)

    # 8) Tests: integrability + orientability
    integrable = path_independence_test(branches)
    orientable = orientation_transport_test(frames)

    return {
        "image_I": σ_e,
        "trace_T": {
            "frames": frames,
            "models": models,
            "branches": branches,
            "remainder": remainder
        },
        "costs": {"C": C, "D": D, "K": K, "demons": demons},
        "tests": {"integrable": integrable, "orientable": orientable},
    }
\end{lstlisting}

\section{Conclusion}

You should now be able to take a small artifact---one sentence, one label, one ``obvious'' statement---and:
\begin{itemize}[leftmargin=2em]
  \item expand it into a locally measurable structure ($\sigma$),
  \item track how different frames impose different apertures and orientations,
  \item see why compression produces distortion even when nobody lies,
  \item test integrability (path dependence) instead of arguing endlessly,
  \item detect orientability failures (sign flips across loops),
  \item locate demons: invariants maintained by exported cost.
\end{itemize}

That is FieldSemantics from first principles: not a worldview, but a disciplined way to keep meaning local, vectorial, and cost-accounted---so systems stop pretending they are coherent by silently billing downstream.

\end{document}